
\documentclass[12pt, a4paper, twoside]{article}

%--Packages--
\usepackage[english]{babel}
\usepackage[utf8]{inputenc}
\usepackage[T1]{fontenc}
\usepackage{showkeys}
\usepackage{latexsym,amsmath,amsthm,amssymb,bm,lmodern,titlesec,xcolor,graphicx,setspace,pdfpages,
geometry,marginnote,caption,array,sectsty,siunitx,eurosym,colortbl,float,wrapfig,vmargin,tikz,hyperref}
\usepackage[nottoc]{tocbibind}
\usepackage{tikz-cd}


\setpapersize{A4}
\setmarginsrb{30mm}{20mm}{30mm}{30mm}%left top right bottom
                 {8mm}{6mm}{10mm}{10mm}%headheight headsep footheight footskip

% Some useful colors
\definecolor{Rosso}{RGB}{232,61,61}
\definecolor{Viola}{RGB}{188,30,188}
\definecolor{Celeste}{RGB}{34,186,211}
\definecolor{Arancione}{RGB}{232,181,19}
\definecolor{Blu}{RGB}{0, 44, 92}
\definecolor{Azzurro}{RGB}{57,186,238}


\hypersetup{colorlinks=true,linkcolor=Blu,allbordercolors=white, urlcolor=blue, citecolor=Viola}
\usepackage{enumitem} 
\renewcommand\labelitemi{\textbullet} 


% math environments
%  we use \newtheorem to specify our theorem-like environments
   % (theorem, definition, etc.) and how to display and number them.
   %
   % Note: The \theoremstyle declarations affect the appearance of the
   % Theorems, Definitions, etc.
   % Note 2: If you want swedish names you have to go and change the argument in brackets e.g. \newtheorem{thm}[defi]{Theorem} should become \newtheorem{thm}[defi]{Sats}
\theoremstyle{definition}
\newtheorem{definition}{Definition}[section]

\newtheorem*{problem}{Problem}
\newtheorem{conjecture}[definition]{Conjecture}

%\theoremstyle{plain}
\newtheorem{theorem}[definition]{Theorem}
\newtheorem{proposition}[definition]{Proposition}
\newtheorem{corollary}[definition]{Corollary}
\newtheorem{lemma}[definition]{Lemma}
%\newtheorem*{remark}{Remark}%[definition]{Remark}
%\newtheorem{example}[definition]{Example}
%\newtheorem*{notation}{Notation}
%\newtheorem*{note}{Note}

\theoremstyle{remark}
\newtheorem*{question}{Question}
\newtheorem{remark}[definition]{Remark}
\newtheorem{example}[definition]{Example}
\newtheorem*{notation}{Notation}
\newtheorem*{note}{Note}

\onehalfspacing
\numberwithin{equation}{section}

%--Heading of  section--
\titleformat{\section}[display]
  {\normalfont\bfseries\color{Blu}}
  {\filleft
    \begin{tikzpicture}
    \node[
      outer sep=0pt,
      text width=2.5cm,
      minimum height=3cm,
      fill=white,
      font=\color{Blu}\fontsize{80}{0}\selectfont,
      align=center
      ] (num) {\thesection};
    \node[
      rotate=90,
      anchor=south,
      font=\color{black}\large\normalfont
      ] at ([xshift=-5pt]num.west) {\Huge\textsc{}};  
    \end{tikzpicture}%
  }
  {10pt}
  {\titlerule\vskip4pt\Huge\sc\sffamily}

%--Heading of subsection--
\newcommand*\numb[1]{%
\begin{tikzpicture}[baseline=-0.7ex]
\node[
  outer sep=0pt,
      text width=0.6cm,
      minimum height=0.6cm,
      fill=Blu,
      font=\color{white}\fontsize{12}{20}\selectfont,
      align=center
      ] (num) {\thesubsection};
\end{tikzpicture}%
}
\titleformat{\subsection}
  {\normalfont\color{Blu}\large\sc\sffamily}{\numb{\thesubsection}}{0.8em}{}

%--Heading subsubsection--
\newcommand*\stocaz[1]{%
\begin{tikzpicture}[baseline=-0.7ex]
\node[
  outer sep=0pt,
      text width=0.8cm,
      minimum height=0.5cm,
      fill=Blu,
      font=\color{white}\fontsize{10}{20}\selectfont,
      align=center
      ] (num) {\thesubsubsection};
\end{tikzpicture}%
}
\titleformat{\subsubsection}
  {\normalfont\color{Blu}\bfseries\sc\sffamily}{\stocaz{\thesubsubsection}}{0.8em}{}
  
% Blackboard Boldface
\newcommand{\NN}{\mathbb{N}}
\newcommand{\ZZ}{\mathbb{Z}}
\newcommand{\QQ}{\mathbb{Q}}
\newcommand{\RR}{\mathbb{R}}
\newcommand{\CC}{\mathbb{C}}
\newcommand{\KK}{\mathbb{K}}
\newcommand{\FF}{\mathbb{F}}

\newcommand{\1}{\mathbf{1}}
\newcommand{\0}{\mathbf{0}}

% Text
\newcommand{\GL}{\text{GL}}
\newcommand{\SL}{\text{SL}}
\newcommand{\SymP}{\text{Sym}}
\newcommand{\AltP}{{\bigwedge}}
\newcommand{\Hom}{\text{Hom}}
\newcommand{\SO}{\text{SO}}
\newcommand{\SU}{\text{SU}}
\newcommand{\id}{\text{id}}
\newcommand{\Stab}{\text{Stab}}
\newcommand{\Cent}{\text{Cent}}
\newcommand{\Norm}{\text{Norm}}
\newcommand{\Type}{\text{Type}}
\newcommand{\Dih}{\text{D}}
\newcommand{\Sym}{\text{S}} %{\mathfrak{S}}
\newcommand{\Alt}{\text{A}}
\newcommand{\Klein}{\text{V}}
\newcommand{\Cyc}{\text{C}}
\newcommand{\sgn}{\text{sgn}}
\newcommand{\Span}{\text{Span}}
\newcommand{\im}{\text{im }}

% Vectors
\newcommand{\bas}{\mathbf{\hat{e}}}

\newcommand{\ihat}{\mathbf{\hat{\imath}}}
\newcommand{\jhat}{\mathbf{\hat{\jmath}}}
\newcommand{\ghat}{\mathbf{\hat{g}}}
\newcommand{\hhat}{\mathbf{\hat{h}}}
\newcommand{\vhat}{\mathbf{\hat{v}}}
\newcommand{\what}{\mathbf{\hat{w}}}
\newcommand{\uhat}{\mathbf{\hat{u}}}

\newcommand{\vvec}{\mathbf{v}}
\newcommand{\uvec}{\mathbf{u}}
\newcommand{\wvec}{\mathbf{w}}


\newcommand{\Ss}{\scriptstyle}



\begin{document}

% This includes the front page (you need to compile the cover.tex separately and have the .pdf available in the same folder as this .tex file).
%\setlength{\voffset}{0cm}\setlength{\hoffset}{0cm}\includepdf[pages={1,2,3,4}]{cover}\setlength{\voffset}{-2.54cm}\setlength{\hoffset}{-2.54cm}



%\clearpage{\thispagestyle{empty}\cleardoublepage}

%\begin{abstract}
%    Your summary goes here
%\end{abstract}

%create the table of contents
\clearpage{\thispagestyle{empty}\cleardoublepage}
{\hypersetup{linkcolor=black}
\tableofcontents
}

%% your first section
\clearpage{\thispagestyle{empty}\cleardoublepage}
\section{Introduction}

Lorem ipsum dolor sit amet, consectetur adipiscing elit, sed do eiusmod tempor incididunt ut labore et dolore magna aliqua.
Ut enim ad minim veniam, quis nostrud exercitation ullamco laboris nisi ut aliquip ex ea commodo consequat. 
Duis aute irure dolor in reprehenderit in voluptate velit esse cillum dolore eu fugiat nulla pariatur. 
Excepteur sint occaecat cupidatat non proident, sunt in culpa qui officia deserunt mollit anim id est laborum.

You can also use swedish letters ä ö å, Ä, Ö, Å. Unicode characters in general should work.


%Another section
\clearpage{\thispagestyle{empty}\cleardoublepage}
\section{Theorems and definitions}

\begin{definition}[Test av definition]
	Test av definition
\end{definition}

\begin{problem}[Test av problem]
	Test av problem
\end{problem}

\begin{conjecture}[Test av conjecture]
	Test av conjecture
\end{conjecture}

\begin{theorem}[Test av theorem]
	Test av theorem
\end{theorem}

\begin{proposition}[Test av proposition]
	Test av proposition
\end{proposition}

\begin{corollary}[Test av corollary]
	Test av corollary
\end{corollary}

\begin{lemma}[Test av lemma]
	Test av lemma
\end{lemma}

\begin{question}[Test av question]
	Test av question
\end{question}

\begin{remark}[Test av remark]
	Test av remark
\end{remark}

\begin{example}[Test av example]
	Test av example
\end{example}

\begin{notation}[Test av notation]
	Test av notation
\end{notation}

\begin{note}[Notering]
	This note's for you
\end{note}


\begin{definition}
	This is a definition!
\end{definition}

\begin{theorem}\label{thm: label}
	This is a theorem!
\end{theorem}

To cross reference: By Theorem~\ref{thm: label} we can say...
\begin{lemma}
	This is a lemma.
\end{lemma}
\begin{proof}
	Put here the proof of the lemma!
\end{proof}


\begin{proof}[Proof of Theorem \ref{thm: label}]
	Another proof
\end{proof}


\begin{corollary}
	here is a corollary!
\end{corollary}

If you want to write a beautiful equation in line do like that $e^{i\pi}+1=0$.
You might also want to write it on its own line
\[
e^{i\pi}+1=0.
\]
If you think that you are going to need to cross reference your equation:
\begin{equation}\label{eq:beautiful}
	e^{i\pi}+1=0
\end{equation}
and this is how you cross reference it \eqref{eq:beautiful}
You might also want to have aligned equations on more lines.
\begin{align}
	e^{i\pi}+1&=0\\
	(\cos\theta)^2+(\sin\theta)^2 &=1
\end{align}

If you need to cross reference something in the bibliography, use the cite command: \cite[p. 45]{Biggs2002}. 
The actual references are stored in a different file, called bibliographyFile.bib 

You can reference most things you found useful,
for example, books \cite{Biggs2002,Knuth1998ArtOfProgramming,LascouxPolynomialsBook}.
Research articles and preprints, \cite{AlexanderssonSulzgruber2019,Sherman2015x}.
Or another thesis, \cite{Teff2013}.

You can also cite online sources. For example, the Encyclopedia of integer sequences, \cite{OEIS},
or a YouTube video \cite{Bazett2018yt}. 
It is even possible to refer to a discussion you had with someone, \cite{PrivateCommunicationWithJustin}.
This is usually rare, it is better to thank them in the acknowledgements section (if they were helpful).



%another section
\clearpage{\thispagestyle{empty}\cleardoublepage}
\section{Tables and graphs}
Here a few examples of tables and graphs.
%A subsection
\subsection{Tables}
\begin{center}
	\begin{tabular}{c c c c c c c c}
		\arrayrulecolor{Azzurro}
		\hline
		{\bfseries Code} & {\bfseries CdL} & {\bfseries xxx} & {\bfseries $T_{setup/lotto}$} & {\bfseries $T_{lav/st}$} & {\bfseries $T_{proc/pezzo}$} & {\bfseries Quantitity} & {\bfseries $T_{tot}$}\\
		\hline
		100 & 4 & 250 & 25 & 0,5 & 0,6 & 1 & 0,6\\
		111 & 2 & 250 & 20 & 2 & 2,08 & 1 & 2,08 \\
		111 & 3 & 250 & 15 & 1,5 & 1,56 & 1 & 1,56 \\
		112 & 2 & 250 & 20 & 2,5 & 2,58 & 1 & 2,58 \\
		112 & 3 & 250 & 15 & 2 & 2,06 & 1 & 2,06\\
		113 & 3 & 500 & 15 & 1 & 1,03 & 2 & 2,06\\
		120 & 1 & 50 & 30 & 2 & 2,6 & 0,1 & 0,26\\
		121 & 1 & 25 & 30 & 3 & 4,2 & 0,1 & 0,42 \\
		121 & 1 & 25 & 30 & 2,5 & 3,7 & 0,1 & 0,37 \\
		\hline
	\end{tabular}
\end{center}



%another section
\clearpage{\thispagestyle{empty}\cleardoublepage}

\section{Insert a figure}
\begin{figure}[H]
	\centering
	\includegraphics[width=1\textwidth]{grafo.png}
	\caption{\label{fig:1} put the caption here}
\end{figure}

\subsection{Footnote}

You can create a footnote like this.\footnote{I created a footnote.}



\clearpage{\thispagestyle{empty}\cleardoublepage}
\section{Conclusion}

Remember the internet is full of places that can help you write a beautiful document with latex.
Some example are \url{overleaf.com} and \verb!stackexchange! just ask Google! 
In addition with an SU e-mail address you should be able to 
create (for free) a premium account on \url{overleaf.com}, 
so you do not need to download a tex editor and you will have
the full latex distribution at your hand (if you are connected to the internet).


$$ \NN \ \ZZ \ \QQ \ \RR \ \CC \ \KK \ \FF$$
$$ \1 \ \0 $$

$$ \GL(V),  \SL_n(\KK), \SymP^k V, \AltP^k V, \Hom(V,W), \SO, \SU, \id, \Stab(g), \Cent(g), \Norm(g),$$ $$ \Type(\sigma), \Dih_{2n}, \Sym_n, \Alt_n, \Klein, \Cyc_n, \sgn(\sigma), \Span(\bas_i)_i $$

$$ \bas \ihat \jhat \ghat \hhat \vhat \what \uhat \vvec \uvec \wvec $$
\clearpage{\thispagestyle{empty}}
\section{Preliminary topics}

\subsection{Linear algebra topics}

	\paragraph{Definition of kernel and image of a map}
	
	\begin{definition}[Kernel and image of a linear map]\label{def:kernelimage}
		Let $V$ and $W$ be two vector spaces and let $\varphi: V \rightarrow W$ be a linear map. Then the kernel and the image of the map are defined thusly:
		\begin{itemize}
			\item[i)] $\ker \varphi = \left\lbrace \vvec \in V \ \middle\vert \ \varphi(\vvec) = \0 \right\rbrace = \varphi^{-1}(\0)$.
			\item[ii)] $\im \varphi = \left\lbrace \wvec \in W \ \middle\vert \ \exists \vvec \in V \text{ s.t. } \varphi(\vvec) = \wvec \right\rbrace = \varphi(V)$.
		\end{itemize}
	\end{definition}
	
	\begin{remark}
		The kernel and the image of a linear map are subspaces of the domain and codomain of the map respectively, ie. $\ker \varphi$ is a subspace of $V$ and $\im \varphi$ is a subspace of $W$.
	\end{remark}
	
	\paragraph{Existence of complementary subspaces}
	
	\begin{corollary}\label{thm:compsubspaces}
		Let $V$ be a vector space and $W$ be a vector subspace of $V$. Then there exists a complementary vector subspace $W'$ in $V$ such that $W \cap W' = \emptyset$ and $W \cup W' = V$. This is equivalent to saying that $V$ is the direct sum of $W$ and $W'$, denoted as $V = W \oplus W'$.
	\end{corollary}
	\begin{proof}
		La la la
	\end{proof}
	
\subsection{Group theoretical topics}

	\subsubsection{Group action}
	
		We could let the group $G$ act on a set $X$ and study the action of $G$ on $X$ defined by
		\begin{align*}
			G \times X \rightarrow X
		\end{align*}
		and thereby permuting the elements of $X$ by for any $g \in G$,
		\begin{align*}
			(g,x) \mapsto g \cdot x = gx,
		\end{align*}
		for any $x \in X$.
	
	\subsubsection{The symmetric group}
	
		Recall that two elements $\sigma, \tau \in \Sym_n$ are conjugate if there exists a $\pi \in \Sym_n$ such that $\sigma = \pi \tau \pi^{-1}$. The conjugacy classes partition $\Sym_n$ into disjoint subsets. 
		
		\begin{proposition}[Conjugacy in $\Sym_n$]
			Any elements $\sigma, \tau \in \Sym_n$ are conjugate if and only if they are of the same cycle type.
		\end{proposition}
		
		\begin{definition}[Sign of a permutation]
			\begin{align*}
				\sgn(\sigma) = (-1)^k,
			\end{align*}
			where $k$ is the number of transpositions required to compose $\sigma$.
		\end{definition}
		
		\begin{proposition}\label{prop:signwelldefined}
			The sign of a permutation is well-defined. That is a permutation is either \textit{even} with sign +1 or \textit{odd} with sign -1.
		\end{proposition}
		\begin{proof}
			\cite[Thm.12.6.1.]{Biggs}
		\end{proof}
		
		
		

\clearpage{\thispagestyle{empty}}
\section{Representations of finite groups}

\paragraph{First definitions}

	Let $G$ be a finite group written multiplicatively. Denote its order by $|G|$ and the action of its elements with $\cdot$, eg. $g \cdot x$.
	
	Let $V$ be a finite-dimensional vector space over the field of complex numbers $\CC$. Its group of automorphisms is denoted by $\GL(V)$. If $V$ is provided with a basis, usually denoted $(\bas_i)_{i=1}^n$, where $n = \dim V$, then $\GL(V)$ is identified with the set of linear transformations of $V$, denoted $\GL_n(\CC)$\cite[18.1]{DummitFoote}.
	
	\begin{definition}[Representation]\cite{Serre}
		Given a finite group $G$, a representation of $G$ in the vector space $V$ is a homomorphism
		\begin{align*}
			\rho: G \rightarrow \GL(V).
		\end{align*}
	\end{definition}
	For every $g$ in $G$, there is a map $\rho_g$ in $\GL(V)$ acting on the elements of $V$ defined by \begin{align*}
		g \cdot \vvec = \rho_g (\vvec), 
	\end{align*}
	and also for any $g,h \in G$ we have
	\begin{align}\label{eq:homomorphism}
		\rho_{gh} = \rho_g \rho_h.
	\end{align}
	From the homomorphism of $\rho$, two consequences follow.
	\begin{proposition}
		For the identity element $e$ of $G$ and an arbitrary element $g$ with inverse $g^{-1}$ in $G$ we have,
		\begin{align*}
			\rho_e = \id \quad \text{and} \quad \rho_g^{-1} = \rho_{g^{-1}},
		\end{align*}
		where $\id$ is the identity transformation and $\rho_{g^{-1}}$ is the inverse of the representation of $g$.
	\end{proposition}
	\begin{proof}
		Take $g$ as an arbitrary element of $G$. The first identity follows from taking $h=e$ in Equation~\ref{eq:homomorphism}:
		\begin{align}
			\rho_{eg} = \rho_e \rho_g.
			%\rho(eg) = \rho(e)\rho(g). %	\rho(g) \overset{!}{=} \rho(eg) = \rho(e)\rho(g)
		\end{align}
		Since $eg=g$ for any $g\in G$ we must have that $\rho_g = \rho_e\rho_g$ which is true if and only if $\rho_e = \id$. Instead taking $h=g^{-1}$ we have:
		\begin{align}
			\rho_{gg^{-1}} = \rho_g \rho_{g^{-1}},
			%\rho(g) \rho(g^{-1}) =  \rho(gg^{-1}~). %\overset{!}{=} \rho(e) = \id.
		\end{align}
		but $gg^{-1} = e$ and $\rho_e = \id$, implying $\rho_{g^{-1}} = \rho_g^{-1}$.
	\end{proof}
	
\paragraph{Choice of basis}
\paragraph{Degree of a representation}

\subsection{Basic examples}
\paragraph{Representations of degree 1, including trivial and alternating}
\paragraph{Permutation representation}
\paragraph{Regular representation}


\subsection{Subrepresentations}
	
\marginnote{Morphisms of reprs.}

	\begin{definition}[$G$-linear map]\cite{FultonHarris}
		Let $V$ and $W$ be two representation spaces of a group $G$. A vector space map 
		\begin{align*}
			\varphi: V \rightarrow W
		\end{align*} 
		is called a $G$-linear map if it commutes with the group action of $G$, ie. for any $\vvec \in V$ and $g \in G$ we have
		\begin{align*}
			\varphi( g \cdot \vvec ) = g \cdot \varphi (\vvec),
		\end{align*}
		or in terms of the maps $\rho_g^V: G \rightarrow \GL(V)$ and $\rho_g^W: G \rightarrow \GL(W)$,
		\begin{align*}
			\varphi \circ \rho^V_g (\vvec) = \rho^W_g \circ \varphi (\vvec).
		\end{align*}
		Equivalently one can say that the diagram in Figure~\ref*{fig:commapmorp} is commutative for every $g \in G$.
		\begin{figure}%[hbt!]
			\centering
			\[\begin{tikzcd}
				V \arrow{r}{\varphi} \arrow[swap]{d}{g \cdot} & W \arrow{d}{g \cdot} \\
				V \arrow{r}{\varphi} & W
			\end{tikzcd}\]
			\caption{Morphism of representations}
			\label{fig:commapmorp}
		\end{figure}
	\end{definition}
	
	
	
	
	
	
	Consider the case where $W$ is a $G$-invariant\footnote{For $W$ to be a $G$-invariant space, it means that for every $g \in G$, $g \cdot \wvec \in W$ for any $\wvec \in W$ or that $\rho_g W \subseteq W$ for any $g\in G$.} vector subspace of $V$% and let $\pi:V \rightarrow W$ be the projection map of $V$ onto $W$. 
	and let $\rho$ be a representation of $G$ in $V$. Then the restriction of $\rho$ to $W$, here denoted $\rho|_W$ is a isomorphism of $W$ onto itself since $W$ is $G$-invariant, and thus $\rho_{gh}|_W = \rho_g|_W \cdot \rho_h|_W$, hence $\rho|_W: G \rightarrow \GL(W)$ is a representation if $G$ in $W$, motivating the following definition.
	
	\begin{definition}[Subrepresentation]\label{def:subrepr}
		Let $G$ be a finite group and $\rho$ be a representation of $G$ in a vector space $V$. A restriction of $\rho$ to a $G$-invariant vector subspace $W$ of $V$ is called a subrepresentation of $\rho$.
	\end{definition}
	
	\begin{proposition}
		Let $V$ and $W$ be representation spaces of a group $G$ and let $\varphi: V \rightarrow W$ be a $G$-linear map. Then the kernel and the image of $\varphi$ are also representations of $G$ and more specifically, $\ker \varphi$ is a subrepresentation of $V$ and $\im \varphi$ is a subrepresentation of $W$.\marginnote{Cokernel??} 
	\end{proposition}
	\begin{proof}
		\begin{itemize}
			\item[i)] Let $\vvec \in \ker \varphi$. Then $g \cdot \varphi (\vvec) = \0$ since by definition $\varphi(\vvec) = \0$, but then we must also have $\varphi ( g \cdot \vvec ) = \0$ since $\varphi$ is a $G$-linear map, implying that $g \cdot \vvec \in \ker \varphi$. Since this holds for any $g \in G$ and $\vvec \in \ker \varphi$, the kernel of $\varphi$ is a $G$-invariant subspace of $V$ and thus a subrepresentation of $V$. \marginnote{Subrepr. not yet defined}
			\item[ii)] For a $\wvec \in \im \varphi$, by definition there exists a $\vvec \in V$ such that $\varphi(\vvec) = \wvec$. Since $g \cdot \vvec \in V$, then we have that $\varphi ( g \cdot \vvec) \in W$. This implies that $g \cdot \varphi (\vvec)$ is also in $W$. Since this holds for any $g \in G$ and $\vvec \in V$, the image of $\varphi$ is a $G$-invariant subspace of $W$ and a thus a subrepresentation of $W$.	\qedhere
%			\item[i)] Let $\vvec \in \ker \varphi$. Then $\rho_g \circ \varphi (\vvec) = \0$, but then we must also have $\varphi \circ \rho_g (\vvec) = \0$, implying that $\rho_g (\vvec) \in \ker \varphi$. Since this holds for any $g \in G$ and any $\vvec \in \ker \varphi$, the kernel of $\varphi$ is a $G$-invariant subspace of $V$ and thus a subrepresentation of $V$.
%			\item[ii)] For a $\wvec \in \im \varphi$, there exists a $\vvec \in V$ such that $\varphi(\vvec) = \wvec$. Since $\rho_g(\vvec) \in V$, then we have that $\varphi \circ \rho_g(\vvec) \in W$. This implies that $\rho_g \circ \varphi (\vvec)$ is also in $W$. Since this holds for any $g \in G$ and $\vvec \in V$, the image of $\varphi$ is a $G$-invariant subspace of $W$ and a thus a subrepresentation of $W$.	\qedhere
		\end{itemize}
	\end{proof}
	
	
	
	
	Recall from linear algebra the existence of complementary vector subspaces of a vector space. Is there a similar property of subrepresentations?
	
	\begin{theorem}[Existence of complementary subrepresentations]
		Let $\rho$ be a representation of a finite group $G$ in a finite-dimensional vector space $V$. Let $W$ be a $G$-invariant subspace of $V$. Then there is a  $G$-invariant subspace $W'$ of $V$ complementary to $W$ such that $V = W \oplus W'$.
	\end{theorem}
	\begin{proof}
		By corollary \ref{thm:compsubspaces}, there is a complementary subspace $W'$ in $V$ such that $V = W \oplus W'$. An arbitrary vector $\vvec \in V$ can be written as a direct sum of a $\wvec$ from $W$ and a $\wvec'$ from $W'$. Let $\pi: V \rightarrow W$ be the orthogonal projection of $V$ onto $W$, then $\pi(\vvec) = \wvec$, $\pi(\wvec) = \wvec$ and $\pi(\wvec') = \0$. 
		
		Consider the ``average'' of this projection over $G$, 
		\begin{align*}
			\bar{\pi} = \frac{1}{|G|} \sum_{g \in G} \rho_g \circ \pi \circ \rho_g^{-1},
		\end{align*}
		and its action on a vector $\wvec \in W$:
		\begin{align*}
			\bar{\pi}(\wvec) = \frac{1}{|G|} \sum_{g \in G} \rho_g \circ \pi \circ \rho_g^{-1} (\wvec).
		\end{align*}
		
		
		
		Since $W$ is $G$-invariant, $\rho_g$ preserves $W$, and so do $\rho_{g}^{-1}$, so $\rho_g^{-1}(\wvec)$ is in $W$. 
		
		By definition, $\pi$ leaves any vector in $W$ fixed, so since $\rho_g^{-1}(\wvec) \in W$, we have $\pi \circ \rho_g^{-1}(\wvec) = \rho_g^{-1}(\wvec)$, so now we have 
		\begin{align*}
			\bar{\pi}(\wvec) = \frac{1}{|G|} \cdot\sum_{g \in G} \underset{=\id}{\underbrace{\rho_g \circ \rho_g^{-1}}} (\wvec)
			= \frac{1}{|G|} \cdot \wvec \sum_{g \in G} 1 = \frac{1}{|G|}\cdot  |G| \cdot \wvec = \wvec,
		\end{align*}
		thus the image of $\bar{\pi}$ is in isomorphic correspondence with the elements of $W$ and also the elements of $V$. Hence $\bar{\pi}$ is a $G$-linear map and behaves like Figure \ref{fig:commapmorp}, that is for any $g \in G$,		\begin{align}\label{eq:glinmap}
			\bar{\pi} \circ \rho_g (\vvec) = \rho_g \circ \bar{\pi} (\vvec)
		\end{align}
		holds for any $\vvec \in V$.
		
		Let us apply Equation~\ref{eq:glinmap} on a $\wvec' \in W'$. Since \
		
		\begin{align*}
			\bar{\pi}(\wvec') = \frac{1}{|G|} \sum_{g \in G} \rho_g \circ \pi \circ \rho_g^{-1}(\wvec')
		\end{align*}
	\end{proof}
	
\subsection{Irreducible representations}
		
		
		
		

%The section with references.

\clearpage{\thispagestyle{empty}}
\bibliographystyle{alphaurl} %This style prints the urls of online sources.


% All references are to be stored in the file bibliographyFile.bib! 
% It is in the same folder as your main .tex file.
% It is using the bibtex syntax, and you can get these entries from google scholar, etc.
% Your supervisor should know,
\bibliography{bibFile}
%\bibliography{bibliographyFile} 

\end{document}
